%============================================================
%  ch02_banach.tex
%  Chapter 2: Banach Space Prerequisites
%
%  Tensor Formats in Banach Spaces
%  Antonio Falcó, CEU Universities
%
%  Estimated length: ~30 pp.
%============================================================

\chapter{Banach Space Prerequisites}
\label{chap:banach}

This chapter collects the functional-analytic and differential-geometric
background required throughout the monograph. The reader familiar with
Banach space theory and infinite-dimensional differential geometry may
proceed directly to Chapter~\ref{chap:algtensor}, referring back as needed.
All vector spaces are over the field $\KK$, which is either $\RR$ or $\CC$;
unless stated otherwise, $\KK = \RR$.

%------------------------------------------------------------
\section{Normed spaces and Banach spaces}
\label{sec:banach:basics}
%------------------------------------------------------------

We begin by fixing notation and recalling the central notions of normed and
Banach space theory that will be used throughout.

\subsection*{Basic definitions}

Let $X$ be a vector space over $\KK$. A \emph{norm} on $X$ is a map
$\|\cdot\|: X \to [0,\infty)$ satisfying:
\begin{enumerate}
\item[(N1)] $\|x\| = 0$ if and only if $x = 0$,
\item[(N2)] $\|\lambda x\| = |\lambda|\,\|x\|$ for all $\lambda \in \KK$,
    $x \in X$,
\item[(N3)] $\|x + y\| \leq \|x\| + \|y\|$ for all $x, y \in X$.
\end{enumerate}
The pair $(X, \|\cdot\|)$ is called a \emph{normed space}. A sequence
$(x_n)_{n \in \NN}$ in $X$ is a \emph{Cauchy sequence} if
$\|x_n - x_m\| \to 0$ as $n, m \to \infty$. A normed space in which every
Cauchy sequence converges is called a \emph{Banach space}.

Two norms $\|\cdot\|$ and $\|\cdot\|'$ on $X$ are \emph{equivalent} if
there exist constants $c, C > 0$ such that
$c\|x\| \leq \|x\|' \leq C\|x\|$ for all $x \in X$.
Equivalent norms define the same topology, and hence the same Cauchy
sequences, so they produce the same Banach space structure.

\begin{example}
\label{ex:banach:lp}
Let $I \subset \RR$ be a measurable set and $1 \leq p < \infty$. The
Lebesgue space $L^p(I)$, equipped with the norm
$\|f\|_{L^p} = \bigl(\int_I |f(x)|^p\,dx\bigr)^{1/p}$,
is a Banach space. For $p = \infty$ one uses the essential supremum norm.
The case $p = 2$ gives a Hilbert space with inner product
$\langle f, g\rangle = \int_I f(x)\overline{g(x)}\,dx$.
\end{example}

\begin{example}
\label{ex:banach:sobolev}
For a bounded open set $\Omega \subset \RR^n$, $1 \leq p < \infty$, and
$k \in \NNz$, the Sobolev space $W^{k,p}(\Omega)$ consists of all functions
$f \in L^p(\Omega)$ whose distributional derivatives $D^\alpha f$ of order
$|\alpha| \leq k$ belong to $L^p(\Omega)$. Equipped with the norm
\begin{equation*}
    \|f\|_{W^{k,p}} = \Bigl(\sum_{|\alpha| \leq k}
        \|D^\alpha f\|_{L^p}^p\Bigr)^{1/p},
\end{equation*}
it is a Banach space. For $p=2$ one writes $H^k(\Omega) = W^{k,2}(\Omega)$,
which is a Hilbert space. These spaces appear as natural component spaces in
high-dimensional tensor problems; see Example~\ref{ex:min:sobolev} and
Section~\ref{sec:topTBT:bestapprox}.
\end{example}

\subsection*{Continuous linear maps}

Let $X$ and $Y$ be normed spaces. A linear map $A: X \to Y$ is
\emph{continuous} if and only if it is \emph{bounded}, meaning
\begin{equation*}
    \opnorm{A}{Y}{X} \coloneqq \sup_{x \neq 0} \frac{\|Ax\|_Y}{\|x\|_X}
    = \sup_{\|x\|_X \leq 1} \|Ax\|_Y < \infty.
\end{equation*}
The quantity $\opnorm{A}{Y}{X}$ is the \emph{operator norm} of $A$.
The space of all continuous linear maps from $X$ to $Y$, denoted
$\cL(X, Y)$, is itself a normed space under the operator norm; it is a
Banach space whenever $Y$ is complete. We write $\cL(X) \coloneqq \cL(X,X)$.
The set of bounded linear automorphisms (bounded invertible maps with
bounded inverse) is
\begin{equation*}
    \GL{X} \coloneqq \{A \in \cL(X) : A \text{ invertible}\}.
\end{equation*}
This is an open subset of $\cL(X)$, and hence a Banach manifold modelled
on $\cL(X)$; see Example~\ref{ex:manifold:GL}.

\subsection*{Dual spaces and reflexivity}

The \emph{topological dual} of a normed space $(X, \|\cdot\|)$ is
$\dual{X} \coloneqq \cL(X, \KK)$, equipped with the dual norm
\begin{equation}
\label{eq:dualnorm}
    \|\varphi\|_{\dual{X}} \coloneqq \sup_{\|x\| \leq 1} |\varphi(x)|.
\end{equation}
The dual $\dual{X}$ is always a Banach space, regardless of whether $X$
is complete. The \emph{bidual} $X^{**} \coloneqq (\dual{X})^*$ contains
$X$ as a canonically embedded subspace via the isometric embedding
$\iota: X \to X^{**}$, $\iota(x)(\varphi) = \varphi(x)$. When $\iota$ is
surjective, $X$ is called \emph{reflexive}. All $L^p$ spaces with
$1 < p < \infty$ are reflexive; $L^1$ and $L^\infty$ are not.

Reflexivity plays an important role in the existence of best approximations.
Recall that a closed set $C$ in a normed space $X$ is \emph{proximinal} if
every point of $X$ has at least one nearest point in $C$. In a reflexive
Banach space, every closed convex set is proximinal
(see, e.g., \cite[p.\,61]{Hackbusch2019}). This fact will be used
extensively in Chapter~\ref{chap:minsubspaces} to prove the existence
of best low-rank tensor approximations.

\subsection*{Convexity and smoothness}

Two geometric properties of the norm are of particular importance for
the Dirac--Frenkel variational principle in Chapter~\ref{chap:DF}.

\begin{definition}
\label{def:strictly_convex}
A normed space $(X, \|\cdot\|)$ is \emph{strictly convex} (or
\emph{rotund}) if $\|x + y\|/2 < 1$ for all $x, y \in X$ with
$\|x\| = \|y\| = 1$ and $x \neq y$.
\end{definition}

\begin{definition}
\label{def:uniformly_convex}
A normed space $(X, \|\cdot\|)$ is \emph{uniformly convex} if for every
$\varepsilon > 0$ there exists $\delta > 0$ such that
$\|x + y\|/2 \leq 1 - \delta$ whenever $\|x\|, \|y\| \leq 1$ and
$\|x - y\| \geq \varepsilon$.
\end{definition}

Every uniformly convex Banach space is reflexive and strictly convex
(Clarkson's theorem). In particular, $L^p$ spaces with $1 < p < \infty$
are uniformly convex. Strict convexity guarantees uniqueness of metric
projections onto closed convex sets, a property needed for the
well-posedness of the Dirac--Frenkel equation.

\begin{definition}
\label{def:smooth}
A Banach space $X$ is \emph{smooth} if the limit
$\lim_{t \to 0} \bigl(\|x + ty\| - \|x\|\bigr)/t$
exists for all $x, y \in X$ with $\|x\| = 1$.
\end{definition}

Smoothness is the dual property to strict convexity in the following
sense: $X$ is smooth if and only if $\dual{X}$ is strictly convex, and
$X$ is strictly convex if and only if $\dual{X}$ is smooth.

%------------------------------------------------------------
\section{Complemented subspaces and projections}
\label{sec:banach:complemented}
%------------------------------------------------------------

The central geometric notion underlying all manifold structures in this
monograph is that of a complemented (or split) subspace. Unlike the
Hilbert space setting, complementability is a non-trivial condition in
general Banach spaces.

\subsection*{Projections}

A bounded linear map $P \in \cL(X)$ is a \emph{projection} if $P^2 = P$.
In that case $\mathrm{Im}\,P$ is a closed subspace, $\mathrm{Ker}\,P$ is
closed, and $X = \mathrm{Im}\,P \oplus \mathrm{Ker}\,P$ (as a direct sum
of closed subspaces). The complementary projection is $\Id{X} - P$.

\begin{definition}
\label{def:split_subspace}
Let $X$ be a Banach space. A closed subspace $U \subset X$ is a
\emph{split subspace} (or \emph{complemented subspace}) of $X$ if
there exists a closed subspace $W \subset X$ such that
$X = U \oplus W$.
The pair $(U, W)$ is then called a \emph{pair of complementary
subspaces}, and $W$ is a \emph{(topological) complement} of $U$ in $X$.
We write $\Proj{U}{W}$ for the unique continuous projection from $X$
onto $U$ along $W$.
\end{definition}

The following proposition characterises split subspaces in terms of
projections and provides three equivalent conditions that will each be
used at different points of the text.

\begin{proposition}[{\cite[Proposition~4.2]{FHN2015}}]
\label{prop:split_equiv}
Let $X$ be a Banach space and $U \subset X$ a closed subspace. The
following conditions are equivalent:
\begin{enumerate}
\item[\textup{(a)}] $U$ is a split subspace of $X$.
\item[\textup{(b)}] There exists $P \in \cL(X)$ with $P^2 = P$ and
    $\mathrm{Im}\,P = U$.
\item[\textup{(c)}] There exists $Q \in \cL(X)$ with $Q^2 = Q$ and
    $\mathrm{Ker}\,Q = U$.
\end{enumerate}
\end{proposition}

In condition~(b), the projection $P$ is none other than $\Proj{U}{W}$
for any choice of closed complement $W$ of $U$; in~(c), one has $Q =
\Proj{W}{U}$.

\begin{proposition}[{\cite[Theorem~4.5]{FHN2015}}]
\label{prop:finitedim_split}
Let $X$ be a Banach space. Every finite-dimensional subspace of $X$ is
a split subspace.
\end{proposition}

In a Hilbert space, every closed subspace is complemented (the orthogonal
complement provides the split). This is no longer true in general Banach
spaces: by a classical result of Lindenstrauss and Tzafriri, a Banach
space in which every closed subspace is complemented is necessarily
isomorphic to a Hilbert space. The failure of complementability in
non-Hilbert spaces is precisely what makes the immersion theorems of
Chapter~\ref{chap:tucker_manifold} non-trivial.

\subsection*{The nilpotent algebra associated to a splitting}

The following result, central to the construction of local charts in
Chapters~\ref{chap:tucker_manifold} and~\ref{chap:TB_manifold},
identifies a particularly well-behaved sub-algebra of $\cL(X)$
associated to any splitting $X = U \oplus W$.

\begin{proposition}
\label{prop:nilpotent_algebra}
Let $X$ be a Banach space and $U, W \in \GrassB{X}$ with $X = U \oplus W$.
Define
\begin{equation*}
    \cL_{(U,W)}(X) \coloneqq
    \bigl\{
        \Proj{W}{U} \circ S \circ \Proj{U}{W} : S \in \cL(X)
    \bigr\}
    \cong \cL(U, W).
\end{equation*}
Then:
\begin{enumerate}
\item[\textup{(i)}] The natural map
    $\cL(U, W) \to \cL_{(U,W)}(X)$,
    $L \mapsto \Proj{W}{U} \circ L \circ \Proj{U}{W}$,
    is an isometric isomorphism.
\item[\textup{(ii)}] $\cL_{(U,W)}(X)$ is a nilpotent sub-algebra of
    $\cL(X)$: for every $L, L' \in \cL_{(U,W)}(X)$ one has $L \circ L' = 0$.
\item[\textup{(iii)}] As a consequence, for every
    $L \in \cL_{(U,W)}(X)$:
    \begin{equation}
    \label{eq:exp_nilpotent}
        \exp(L) \coloneqq \sum_{n=0}^{\infty} \frac{L^n}{n!}
        = \Id{X} + L,
        \qquad
        \exp(-L) = \Id{X} - L = (\Id{X} + L)^{-1}.
    \end{equation}
\end{enumerate}
\end{proposition}

\begin{proof}
For the isometry in~(i), note that
$\|\Proj{W}{U} \circ L \circ \Proj{U}{W}\|_{X \leftarrow X}
= \|L\|_{W \leftarrow U}$ by the definition of the operator norm.
For~(ii): if $L = \Proj{W}{U} \circ S \circ \Proj{U}{W}$ and
$L' = \Proj{W}{U} \circ S' \circ \Proj{U}{W}$, then
\begin{equation*}
    L \circ L' =
    \Proj{W}{U} \circ S \circ (\Proj{U}{W} \circ \Proj{W}{U})
    \circ S' \circ \Proj{U}{W} = 0,
\end{equation*}
since $\Proj{U}{W} \circ \Proj{W}{U} = 0$. Statement~(iii) follows
immediately from~(ii) and the series definition of the exponential.
\end{proof}

The identity $\exp(L) = \Id{X} + L$ means that the exponential map on
the nilpotent algebra $\cL_{(U,W)}(X)$ is a polynomial of degree one,
and in particular a global analytic diffeomorphism. This is the key
that makes the local charts of the Grassmann--Banach manifold and the
Tucker manifold globally well-defined; see
Section~\ref{sec:banach:grassmann}.

\subsection*{Common complements and chart compatibility}

The following lemma, which underlies the compatibility of overlapping
charts in the Grassmann--Banach manifold, characterises when two split
subspaces share a common complement.

\begin{lemma}[{\cite[Lemma~2.1]{DixmierDouady1963}}]
\label{lem:common_complement}
Let $X$ be a Banach space and let $U, U', W \in \GrassB{X}$ with
$X = U \oplus W = U' \oplus W$. The following conditions are equivalent:
\begin{enumerate}
\item[\textup{(a)}] $U$ and $U'$ have $W$ as a common complement, i.e.,
    $X = U \oplus W = U' \oplus W$.
\item[\textup{(b)}] The restriction $\Proj{U}{W}\big|_{U'} : U' \to U$
    is a linear isomorphism (with bounded inverse).
\end{enumerate}
Moreover, in this case
$\bigl(\Proj{U}{W}\big|_{U'}\bigr)^{-1} = \Proj{U'}{W}\big|_U$.
\end{lemma}

%------------------------------------------------------------
\section{The Grassmann--Banach manifold}
\label{sec:banach:grassmann}
%------------------------------------------------------------

The Grassmannian of a Banach space is the natural parameter space for
subspaces, and it serves as the base manifold for the Tucker and
tree-based tensor manifolds.

\begin{definition}
\label{def:grassmannian}
Let $X$ be a Banach space. The \emph{Grassmann manifold} (or
\emph{Grassmannian}) of $X$ is
\begin{equation*}
    \GrassB{X} \coloneqq \{U \subset X : U \text{ is a split subspace of } X\}.
\end{equation*}
For $r \in \NNz$, the \emph{finite Grassmannian} is
\begin{equation*}
    \Grass{r}{X} \coloneqq
    \{U \in \GrassB{X} : \dim U = r\}.
\end{equation*}
\end{definition}

Both $X$ and $\{0\}$ belong to $\GrassB{X}$.
By Proposition~\ref{prop:finitedim_split}, every $r$-dimensional subspace
of $X$ belongs to $\GrassB{X}$, so $\Grass{r}{X}$ is well defined.

\begin{theorem}[Grassmann--Banach manifold, {\cite{Douady1966}}]
\label{thm:grassmann_manifold}
Let $X$ be a Banach space. The Grassmannian $\GrassB{X}$ is an analytic
Banach manifold, with atlas
$\bigl\{\GrassB{W, X},\, \Psi_{U \oplus W}\bigr\}_{U \in \GrassB{X}}$,
where for each splitting $X = U \oplus W$:
\begin{enumerate}
\item[\textup{(i)}] $\GrassB{W, X} \coloneqq
    \{V \in \GrassB{X} : X = V \oplus W\}$
    is the chart domain at $U$.
\item[\textup{(ii)}] The chart map
    $\Psi_{U \oplus W}: \GrassB{W, X} \to \cL(U, W)$
    is defined by
    \begin{equation}
    \label{eq:grassmann_chart}
        \Psi_{U \oplus W}(V) \coloneqq
        \Proj{W}{U}\big|_V \circ \bigl(\Proj{U}{W}\big|_V\bigr)^{-1},
    \end{equation}
    with inverse
    $\Psi_{U \oplus W}^{-1}(L) = G(L) \coloneqq
    \{(\Id{X} + L)(u) : u \in U\}$.
\end{enumerate}
The transition maps are analytic. Each chart domain
$\GrassB{W, X}$ is a Banach manifold modelled on $\cL(U, W)$.
\end{theorem}

The chart map $\Psi_{U \oplus W}$ sends a subspace $V \in \GrassB{W,X}$
to the linear map $L: U \to W$ whose graph $G(L)$ equals $V$. The
identity $\Psi_{U\oplus W}^{-1}(L) = \exp(L)(U)$, combined with
Proposition~\ref{prop:nilpotent_algebra}, shows that the inverse chart is
a degree-one polynomial map (hence analytic). Analyticity of the
transition maps follows by a direct computation using
Lemma~\ref{lem:common_complement}.

\begin{remark}
\label{rem:grassmann_not_modelled}
The full Grassmannian $\GrassB{X}$ is a Banach manifold but is
\emph{not} modelled on a single Banach space: if $U$ and $U'$ are
finite-dimensional with $\dim U \neq \dim U'$, then $\cL(U, W)$ and
$\cL(U', W')$ are non-isomorphic. This is not a deficiency but a
structural fact: $\GrassB{X}$ decomposes into connected components
$\GrassB{X} = \bigsqcup_{r \in \NNz} \Grass{r}{X}$, each of which
\emph{is} modelled on a single Banach space.
\end{remark}

\begin{proposition}
\label{prop:grassr_connected}
For each $r \in \NNz$, the finite Grassmannian $\Grass{r}{X}$ is a
connected component of $\GrassB{X}$, and it is an analytic Banach manifold
modelled on $\cL(U, W)$ for any $U \in \Grass{r}{X}$ and complement $W$
\textup{(}the isomorphism class of $\cL(U, W)$ does not depend on the
choice of $U$ or $W$\textup{)}.
\end{proposition}

\begin{remark}[Grassmannian of a normed space]
\label{rem:grassmann_normed}
The construction extends to normed (non-complete) spaces. Let
$(X, \|\cdot\|)$ be a normed space with completion $\bar X$. Define
$\Grass{r}{X}$ as the set of $r$-dimensional subspaces of $X$ that are
also split in $\bar X$. By Lemma~\ref{lem:normed_grassmann}, every chart
domain $\GrassB{W, \bar X}$ is contained in $\Grass{r}{X}$, so the atlas
of $\Grass{r}{\bar X}$ restricts to an analytic atlas on $\Grass{r}{X}$.
This variant is needed in Section~\ref{sec:tucker_manifold:charts}, where
the component spaces $V_\alpha$ are normed but not necessarily complete.
\end{remark}

The tangent space to $\GrassB{X}$ at $U$ is identified as follows.

\begin{proposition}
\label{prop:grassmann_tangent}
For $U \in \GrassB{X}$ with complement $W$, the tangent space
$T_U \GrassB{X} \cong \cL(U, W)$.
\end{proposition}

This identification is consistent with the finite-dimensional case: in
$\RR^n$, the tangent space to the Grassmannian $Gr(r, n)$ at a subspace
$U$ is the space of linear maps from $U$ to its orthogonal complement,
recovering the classical description.

%------------------------------------------------------------
\section{Banach manifolds}
\label{sec:banach:manifolds}
%------------------------------------------------------------

We recall the definition of a Banach manifold in its general form, since
the full Grassmannian $\GrassB{X}$ and the full tensor space
$\VDnorm$ are manifolds not modelled on a fixed Banach space.

\begin{definition}
\label{def:banach_manifold}
Let $\mathbb{M}$ be a set. An \emph{atlas of class $\mathcal{C}^p$}
(for $p \geq 0$ or $p = \infty$ or $p = \omega$, the analytic case)
on $\mathbb{M}$ is a family of \emph{charts}
$\{(M_i, u_i)\}_{i \in A}$ satisfying:
\begin{enumerate}
\item[\textup{(AT1)}] $\{M_i\}_{i\in A}$ covers $\mathbb{M}$.
\item[\textup{(AT2)}] For each $i \in A$, the chart $u_i: M_i \to U_i$
    is a bijection onto an open set $U_i$ in some Banach space $X_i$,
    and for any $i, j \in A$, the set $u_i(M_i \cap M_j)$ is open in $X_i$.
\item[\textup{(AT3)}] Whenever $M_i \cap M_j \neq \emptyset$,
    the \emph{transition map}
    $u_j \circ u_i^{-1} : u_i(M_i \cap M_j) \to u_j(M_i \cap M_j)$
    is a diffeomorphism of class $\mathcal{C}^p$.
\end{enumerate}
Two atlases are \emph{compatible} if their union is again an atlas. A
\emph{$\mathcal{C}^p$-Banach manifold} is a set $\mathbb{M}$ together
with a maximal (complete) atlas of class $\mathcal{C}^p$. If all
$X_i$ are isomorphic to a single Banach space $X$, we say $\mathbb{M}$
is \emph{modelled on $X$}.
\end{definition}

\begin{example}
\label{ex:manifold:banach_space}
Every Banach space $X$ is a $\mathcal{C}^\infty$-Banach manifold
modelled on itself, with the single chart $(X, \Id{X})$.
\end{example}

\begin{example}
\label{ex:manifold:GL}
If $X$ is a Banach space, then $\GL{X}$ is an open subset of $\cL(X)$,
and hence an analytic Banach manifold modelled on $\cL(X)$. The maps
$A \mapsto A^{-1}$ and $(A,B) \mapsto AB$ are analytic
\textup{(}see~\cite[2.7, 2.8]{Upmeier1985}\textup{)}.
\end{example}

\begin{definition}
\label{def:tangent_space}
Let $\mathbb{M}$ be a $\mathcal{C}^p$-Banach manifold with $p \geq 1$
and $m \in \mathbb{M}$. A \emph{tangent vector} at $m$ is an equivalence
class of triples $(U_i, u_i, v)$ where $(U_i, u_i)$ is a chart at $m$
and $v \in X_i$, under the equivalence
$(U_i, u_i, v) \sim (U_j, u_j, w)$ if $D(u_j \circ u_i^{-1})(u_i(m)) v = w$.
The \emph{tangent space} $T_m \mathbb{M}$ is the set of all tangent vectors
at $m$, which inherits the structure of a Banach space from any chart.
\end{definition}

\begin{definition}
\label{def:immersion}
A $\mathcal{C}^1$-map $f: \mathbb{M} \to \mathbb{N}$ between Banach
manifolds is an \emph{immersion} at $m$ if $Df(m): T_m\mathbb{M} \to
T_{f(m)}\mathbb{N}$ is injective with closed split image. It is an
\emph{immersion} if it is an immersion at every point. An
\emph{immersed submanifold} of $\mathbb{N}$ is the image of an injective
immersion.
\end{definition}

The closedness and split condition on the image of $Df(m)$ is automatic
in finite dimensions and in Hilbert spaces (where every closed subspace is
split), but is a genuine constraint in Banach spaces. This is the key
difficulty resolved in Section~\ref{sec:tucker_manifold:immersion}.

%------------------------------------------------------------
\section{Multilinear maps and Fréchet differentiability}
\label{sec:banach:multilinear}
%------------------------------------------------------------

Since the tensor product map is multilinear, the following general results
will be used repeatedly.

\begin{definition}
A map $M: X_1 \times \cdots \times X_d \to Y$ between normed spaces is
\emph{multilinear} if it is linear in each argument separately. It is
\emph{bounded} (equivalently, \emph{continuous}) if
\begin{equation*}
    \|M\| \coloneqq \sup_{\substack{x_j \in X_j \\ \|x_j\|_j \leq 1}}
    \|M(x_1, \ldots, x_d)\|_Y < \infty.
\end{equation*}
\end{definition}

The following proposition shows that continuity implies a much stronger
regularity.

\begin{proposition}[{\cite[Proposition~79]{HJ2012}}]
\label{prop:multilinear_smooth}
Let $X_1, \ldots, X_d$ and $Y$ be normed spaces and let
$M: X_1 \times \cdots \times X_d \to Y$ be a continuous multilinear map.
Then $M$ is $\mathcal{C}^\infty$-Fréchet differentiable, and
$D^k M = 0$ for all $k > d$.
\end{proposition}

The vanishing of derivatives beyond order $d$ reflects the polynomial
nature of multilinear maps. In the complex case, a stronger conclusion
holds.

\begin{proposition}[{\cite[Theorem~160]{HJ2012}}]
\label{prop:analytic_c1}
Let $X$, $Y$ be complex Banach spaces and $U \subset X$ open. A map
$f: U \to Y$ is analytic if and only if it is $\mathcal{C}^1$-Fréchet
differentiable.
\end{proposition}

\begin{corollary}
\label{cor:multilinear_analytic}
Every continuous multilinear map between complex Banach spaces is analytic.
\end{corollary}

In real Banach spaces, $\mathcal{C}^\infty$ and analytic are distinct
notions; the tensor product map is $\mathcal{C}^\infty$ but need not be
analytic over $\RR$. For the real case, $\mathcal{C}^\infty$-regularity
is sufficient for the manifold constructions in Chapters~\ref{chap:tucker_manifold}
and~\ref{chap:TB_manifold}.

\begin{remark}
\label{rem:regularity_manifold}
Proposition~\ref{prop:multilinear_smooth} is the reason why the Tucker
and tree-based tensor manifolds are $\mathcal{C}^\infty$-Banach manifolds
(and analytic over $\CC$): their chart transition maps are compositions of
continuous multilinear maps with bounded linear projections, all of which
are $\mathcal{C}^\infty$-Fréchet differentiable.
\end{remark}

\subsection*{The Grassmannian revisited: Lemma on normed spaces}

For later use in Section~\ref{sec:tucker_manifold:charts}, we record
the following lemma about Grassmannians of normed (non-complete) spaces.

\begin{lemma}
\label{lem:normed_grassmann}
Let $(X, \|\cdot\|)$ be a normed space with completion $\bar X$. Let
$U \in \Grass{r}{\bar X}$ with $U \subset X$, and suppose
$\bar X = U \oplus W$ for some $W \in \GrassB{\bar X}$. Then every
$V \in \GrassB{W, \bar X}$ satisfies $V \subset X$.
\end{lemma}

\begin{proof}
Observe that $X = U \oplus (W \cap X)$, where $W \cap X$ is dense in $W$.
Suppose for contradiction that there exists $V \in \GrassB{W, \bar X}$
with $V \not\subset X$. Since $V$ is finite-dimensional and $\bar X =
V \oplus W$, the subspace $V \cap X$ is a proper subspace of $V$. On
the other hand, $X = (V \cap X) \oplus (W \cap X)$ implies $\bar X =
(V \cap X) \oplus W$, so $V \cap X \in \GrassB{W, \bar X}$. But
$\dim(V \cap X) < \dim V = r$, contradicting the fact that every element
of $\GrassB{W, \bar X}$ has dimension $r$.
\end{proof}


%------------------------------------------------------------
\section{Historical notes}
\label{sec:banach:history}
%------------------------------------------------------------

The theory of Banach spaces was developed primarily in the 1920s and
1930s, with Stefan Banach's monograph~\cite{Banach1932} providing the
first systematic treatment. The duality theory and reflexivity go back
to Banach and Alaoglu; the theorem that uniform convexity implies
reflexivity is due to Clarkson (1936) and Milman (1938).

\paragraph{Complemented subspaces.}
The study of complemented subspaces has its roots in the work of
Banach himself, who showed that $c_0$ has uncomplemented subspaces.
The decisive result --- that a Banach space in which every closed
subspace is complemented must be isomorphic to a Hilbert space ---
was proved by Lindenstrauss and Tzafriri in 1971. This explains why
complementability is a genuine and non-trivial constraint in the
Banach tensor setting of this monograph.

The characterisation of split subspaces via bounded projections
(Proposition~\ref{prop:split_equiv}) and the finite-dimensionality
criterion (Proposition~\ref{prop:finitedim_split}) are classical; the
formulation used here follows~\cite{FHN2015}. The nilpotent algebra
associated to a splitting (Proposition~\ref{prop:nilpotent_algebra})
and its use to identify the chart inverse as a degree-one polynomial
map is a key observation of~\cite{FHN2019}.

The lemma on common complements (Lemma~\ref{lem:common_complement})
is a version of results going back to Dixmier and Douady~\cite{DixmierDouady1963},
who studied continuous fields of Hilbert spaces. The formulation for
general Banach spaces and its application to the transition maps of the
Tucker manifold appears in~\cite{FHN2019}.

\paragraph{Grassmann--Banach manifold.}
The infinite-dimensional Grassmannian was constructed as an analytic
Banach manifold by Douady~\cite{Douady1966} in the context of moduli
spaces of compact analytic subspaces. The atlas described in
Theorem~\ref{thm:grassmann_manifold} goes back to this work; the
specific formulation in terms of the chart map~\eqref{eq:grassmann_chart}
and the identification of the tangent space with $\cL(U,W)$ follows
the presentation of~\cite{Upmeier1985}. The finite-dimensional
Grassmannian $Gr(r,n)$ is a smooth compact manifold of dimension
$r(n-r)$, a fact well known since Grassmann's original work in the
1840s; the infinite-dimensional version requires the split condition
as a substitute for compactness arguments.

\paragraph{Multilinear maps and Fréchet differentiability.}
The fact that continuous multilinear maps are $\mathcal{C}^\infty$
(Proposition~\ref{prop:multilinear_smooth}) is a standard result of
infinite-dimensional calculus; see~\cite{HJ2012} for a careful
treatment. The equivalence between $\mathcal{C}^1$ and analyticity
in complex Banach spaces (Proposition~\ref{prop:analytic_c1}) is the
Gâteaux--Cauchy theorem; see~\cite{HJ2012}. These two propositions are
the main tools that allow one to conclude that the Tucker and tree-based
tensor manifolds are $\mathcal{C}^\infty$ (real case) or analytic
(complex case) without further calculation: their atlas transition maps
are compositions of continuous multilinear maps and bounded linear
projections, all of which satisfy the hypotheses of
Propositions~\ref{prop:multilinear_smooth} and~\ref{prop:analytic_c1}.
